\chapter{Introduction}




In 1937, the Italian physicist Ettore Majorana showed that the relativistic wave equation for fermions admits a completely real representation. In his paper \textit{Teoria simmetrica dell'elettrone e del positrone} (\textit{Symmetric theory of the electron and positron}), he described the possibility of a neutral fermion that is identical to its own antiparticle \cite{majoranaTeoriaSimmetricaDellelettrone1937,narozhnyMajoranaFermionsNonuniform2017}. No elementary Majorana fermion has yet been conclusively identified. However, collective excitations in solid-state systems can behave as Majorana quasiparticles, providing a possible route for studying Majorana physics in controlled devices \cite{mourikSignaturesMajoranaFermions2012, brouwerEnterMajoranaFermion2012,MajoranaFeiMiZiYuTuoBuLiangZiJiSuan2017}.

The interest in these quasiparticles is closely connected to one of the central challenges of quantum computation: quantum information is highly sensitive to environmental noise. Fluctuations in temperature, electromagnetic fields, and imperfect control can cause decoherence and introduce errors into a calculation \cite{ayukaryanaQuestHopeMajorana2021}. Majorana zero modes offer a possible way to reduce sensitivity to some local disturbances. An ordinary fermionic mode can be decomposed into two Majorana components, and in a one-dimensional topological superconductor these components may become spatially separated at opposite ends of the system \cite{aguadoNewTrendsPlatforms2024,ayukaryanaQuestHopeMajorana2021, harleObservingBraidingTopological2023}. The corresponding fermionic degree of freedom is encoded non-locally. Local perturbations then have limited access to the complete encoded state, provided that fermion parity is preserved, the excitation gap remains open, and overlap between the two Majorana components is sufficiently small \cite{ayukaryanaQuestHopeMajorana2021,castagnoliNotionsSymmetryComputational1993}.

Majorana zero modes are also predicted to exhibit non-Abelian exchange statistics. Adiabatically exchanging two modes transforms the state within a degenerate subspace and can implement a quantum gate \cite{lianTopologicalQuantumComputation2018,MajoranaFeiMiZiYuTuoBuLiangZiJiSuan2017}. Such an exchange may be realized by physically moving the modes or by tuning couplings in a sequence that produces the same exchange operation. The possibility of combining non-local encoding with exchange-based gates is what makes Majorana systems interesting for topological quantum computation.

The minimal theoretical model illustrating this physics is the Kitaev chain, introduced by Alexei Kitaev in 2000 \cite{kitaevUnpairedMajoranaFermions2001}. It describes a one-dimensional chain of spinless fermions with nearest-neighbor hopping and \(p\)-wave superconducting pairing. In its topological phase, the model supports Majorana zero modes at its ends and provides a clear description of how non-local fermionic states and exchange operations can emerge. The Kitaev chain is highly idealized, and realizing its ingredients in a controllable physical system remains challenging \cite{sanchesRevisitingPoorMans2026}.

One compact engineered platform inspired by the Kitaev chain is the quantum-dot--superconductor setup, often referred to as a Poor Man's Majorana (PMM) system \cite{sanchesRevisitingPoorMans2026,leijnseParityQubitsPoor2012}. Quantum dots act as discrete sites, while superconducting segments generate effective hopping and non-local pairing through elastic cotunnelling and crossed Andreev reflection. The dot energies, pairing strengths, tunnel couplings, and interaction strengths can be tuned, making PMM systems useful for exploring Majorana-like physics in finite devices. Their finite size is also a fundamental limitation, because the resulting modes are not protected in the same way as the edge modes of an infinite topological system, and low-energy degeneracy alone does not establish that a configuration has useful Majorana character or supports the desired exchange operation \cite{zhangPoorMansMajorana2026}.

The analytical non-interacting two-dot sweet spot is known, but longer interacting chains contain many tunable parameters and no simple general sweet-spot condition \cite{svenssonQuantumDotBased2024}. Numerical optimization offers a way to search this parameter space, but optimizing energy degeneracy alone could also favor ordinary accidental degeneracies. A useful search must therefore combine several complementary criteria, including even--odd parity degeneracy, excitation gaps, local charge indistinguishability, and edge-localized Majorana polarization. Finding a favorable static configuration is still not sufficient: the stronger operational test is whether it supports the expected Majorana exchange under a braiding protocol. This thesis additionally asks whether the effective exchange structure remains meaningful when the projected space is enlarged beyond the lowest manifold.

The investigation is organized around four research questions:
\begin{enumerate}
  \item Can numerical optimization recover known PMM sweet spots and identify promising Majorana-like configurations in interacting two-, three-, and four-dot chains?
  \item Which configurations best balance parity degeneracy, excitation gaps, local charge indistinguishability, and edge-localized Majorana polarization?
  \item Can projected microscopic junction couplings reproduce the ideal Majorana exchange within a selected low-energy manifold?
  \item How do energy-level and projected local-operator braid constructions behave across higher energy sectors, and what does the internal block structure of these sectors imply for logical gate operations?
\end{enumerate}

\section{Scope and Interpretation}

The systems studied in this thesis are finite, spinless PMM chains represented in the full many-body Fock basis and analyzed by exact diagonalization. Throughout the thesis, the term ``Majorana-like'' refers to finite-system configurations that satisfy the selected spectral and local diagnostics. It does not imply exact topological protection or prove the existence of perfectly isolated Majorana zero modes, but that is why the PMM system is interesting: it can be tuned to produce low-energy states that are close enough to the Majorana ideal to make the physical interpretation meaningful and the braiding behavior nontrivial. The numerical optimization identifies favorable minima of the chosen loss function, but it cannot guarantee that the global optimum has been found.

The braiding calculations are also performed within controlled projected models. In the sector-resolved scans, the energy-level, direct local-operator, and block-matched local-operator constructions are each compared with a target defined from their own exchange operators. These self-target comparisons test whether each construction executes its intended braid and how that performance depends on the selected energy sector. They do not by themselves establish that the local constructions reproduce the same exchange channel as the energy-level construction. That stronger statement requires a common-target comparison.

\section{Method Overview}

The first part of the numerical study searches for favorable PMM chains. The hopping amplitudes are fixed to set the energy scale, while the onsite energies and pairing amplitudes are optimized for several chain lengths and fixed interaction strengths. The loss function combines parity-pair splitting, excitation gaps, local charge differences, and Majorana-polarization profiles. Local gradient-based optimization is combined with repeated restarts and basin hopping to explore the resulting non-convex parameter landscape. The selected configurations are then reconstructed and evaluated using parity-resolved spectra and local Majorana diagnostics.

The second part tests braiding. The strongest candidate of the interacting configurations is used as the building block for a three-subsystem junction. Within its lowest projected manifold, an ideal four-Majorana braid is compared with braids generated by projected microscopic junction operators and projected local Majorana components. The analysis is then extended sector by sector using energy-level Majoranas, direct projected local operators, and block-matched local operators. For each construction, the geometric unitary generated by the braid path is computed using Kato transport and compared with the corresponding target exchange.


\section{Thesis Outline}

Chapter 2 presents the theoretical background, including Majorana zero modes, the Kitaev chain and PMM systems, diagnostics of Majorana-like states, Majorana exchange, projected subspaces, and Kato transport. Chapter 3 describes the numerical PMM representation, optimization strategy, post-optimization diagnostics, effective Majorana construction, braiding models, and braid-validation procedure. Chapter 4 presents the optimized configurations, selects the strongest braiding candidate, benchmarks the braid in the lowest projected manifold, and extends the analysis across the identified energy sectors. Chapter 5 discusses the physical interpretation, quantum-computing implications, and limitations of the results. Chapter 6 summarizes the principal conclusions and outlines directions for future work.
% Quantum computing offers the possibility of solving classes of problems that are intractable on classical hardware, but this promise is limited by one central difficulty: quantum states are fragile. Decoherence, noise, and imperfect control make it challenging to preserve and manipulate quantum information long enough to perform useful computations. One of the most attractive proposed solutions to this problem is topological quantum computing, where information is stored non-locally in collective quantum states that are intrinsically less sensitive to local perturbations. In this picture, protection is not achieved only through better engineering, but through the structure of the quantum states themselves.

% Majorana zero modes have therefore attracted broad interest as candidate building blocks for fault-tolerant quantum computing. In condensed-matter systems, they appear not as fundamental particles, but as emergent zero-energy quasiparticles in topological superconductors \cite{Topocondmat}. Their most important feature is that a pair of spatially separated Majorana modes can encode a non-local fermionic degree of freedom. Because local perturbations cannot easily access this information, such states offer a route toward more robust quantum operations. In addition, Majorana modes are predicted to obey non-Abelian exchange statistics, meaning that exchanging them can implement quantum gates through braiding rather than through purely local control pulses.

% Although the basic theoretical picture is well established, realizing and controlling Majorana modes in realistic devices remains a major challenge. Much of the experimental and theoretical effort has therefore focused on engineered hybrid systems, where conventional superconductivity, tunneling, spin structure, and confinement can be combined to produce effective topological behavior \cite{Qtech}. Among the simplest of these platforms is the quantum-dot--superconductor--quantum-dot setup, often referred to as the Poor Man's Majorana system. This device can be understood as a minimal analogue of the Kitaev chain: hopping between dots plays the role of the kinetic term, non-local pairing induced by the superconductor plays the role of effective $p$-wave pairing, and gate-controlled dot energies act as local chemical potentials.

% The appeal of this platform is twofold. First, it is conceptually simple enough to connect directly to the Kitaev-chain picture, which makes it well suited for theoretical analysis. Second, it is experimentally tunable: dot energies, effective pairing strengths, and tunnel couplings can all be adjusted in a controlled way. At the same time, this simplicity comes with limitations. In finite and interacting systems, one should be careful not to overstate the meaning of low-energy degeneracies or Majorana-like observables. The relevant question is therefore not only whether a model displays signatures associated with Majorana physics, but whether those signatures survive under interactions and whether they remain useful for braiding operations.

% This thesis studies that question in interacting chains of coupled quantum dots. More specifically, I investigate whether an interacting $n$-dot Poor Man's Majorana chain can be numerically optimized to produce low-energy states with the expected Majorana-like characteristics, and whether those optimized states support a useful braiding protocol. The goal is not to claim exact topological protection in a finite microscopic model. Rather, the aim is to identify parameter regimes where the low-energy structure is sufficiently close to the Majorana ideal to make the physical interpretation meaningful and the braiding behavior nontrivial.

% To do this, I model the system with an interacting quantum-dot Hamiltonian containing on-site energies, nearest-neighbor hopping, induced pairing, and Coulomb repulsion. The search for promising parameter regimes is formulated as an optimization problem, where the loss function is built from several physically motivated Majorana criteria: near-degeneracy between parity sectors, finite low-energy gaps, small charge differences between parity partners, and a Majorana-polarization profile consistent with edge-localized modes. Once optimized configurations are found, I analyze their parity-resolved spectra and local diagnostics, reconstruct effective Majorana operators from the low-energy manifold, and compare an ideal four-Majorana braiding protocol with a braid generated by projected microscopic junction operators.

% The main results show a clear distinction between weakly and strongly interacting regimes. The optimization reproduces the known non-interacting two-dot sweet spot and, for longer chains, finds more inhomogeneous parameter profiles with strong central structure in the on-site energies and pairing amplitudes. Among the interacting cases studied here, the three-dot system at weak interaction emerges as the most convincing candidate: it combines a very small ground-state splitting, a clear gap to excited states, and a spatial Majorana profile concentrated at the chain edges. In the lowest projected manifold, the corresponding microscopic braid agrees remarkably well with the ideal effective-model braid. However, this agreement deteriorates as stronger interactions and higher-energy sectors are included, indicating that useful braiding depends sensitively on how well the exchange remains confined to a clean low-energy subspace.

% The broader motivation of the thesis is therefore not only to identify candidate Majorana-like configurations, but also to clarify the boundary between effective low-energy success and microscopic breakdown. In that sense, the work is positioned between abstract Majorana models and realistic device-level questions: the optimization tells us which interacting parameter regimes are promising, while the braiding analysis tests whether those regimes support the operation that ultimately matters for topological quantum computing.

% The remainder of the thesis is organized as follows. Chapter 2 presents the theoretical background, beginning with superconductivity and the Bogoliubov-de Gennes formalism, then introducing the Kitaev chain, the quantum-dot realization of Poor Man's Majoranas, and the basic logic of braiding. Chapter 3 describes the numerical methods used to optimize the interacting dot chains, analyze the resulting low-energy states, and construct the braiding protocol. Chapter 4 presents the optimization results, local Majorana diagnostics, and braiding benchmarks in both the idealized and projected microscopic models. Chapter 5 discusses the physical interpretation of these findings, their limitations, and the role of interactions in shaping the low-energy manifold. Finally, Chapter 6 summarizes the main conclusions and outlines possible directions for future work.
